\subsection{Datasets}

In this work, we evaluated our method on two disease classification datasets ISIC 2019 challenge \cite{Combalia2019BCN20000DL} \cite{Tschandl2018TheHD} and the Fitzpatrick-17k dataset \cite{Groh2021EvaluatingDN} . ISIC 2019 challenge contains 25,331 images in 9 diagnostic categories of target labels. We took the gender labels as our sensitive attribute.
The Fitzpatrick-17k dataset contains 16,577 images in 114 skin conditions of target labels and defines the skin tone labels as the sensitive attributes. The skin tones are categorized into six levels from 1 to 6. The smaller the value, the lighter the skin. Therefore, we categorize levels 1 to 3 as light skin and 4 to 6 as dark skin. 

\subsection{Evaluation Metrics}
\label{sec:evaluation_metrics}
Several fairness metrics are proposed to evaluate the degree of fairness in classification task. Demographic parity \cite{dwork2012fairness} and equalized odds (EO) \cite{hardt2016equality,dwork2012fairness} are two well-known fairness criterion. We first define an input feature $X \in \mathbb{R}^d$ with sensitive attribute $A$, ground truth target class $Y$, and the predicted target class $\hat{Y}$. Demographic parity is satisfied if 
\begin{equation}
\label{eq.DEO}
P(\hat{Y} = 1|A = 0) = P(\hat{Y} = 1|A = 1).
\end{equation}
The drawback of demographic parity is that the classifier could achieve the fairness condition by adjusting the proportion of correct rate of two sensitive attributes through misclassifying some instances. On the other hand, EO forces the true positive rate and the false positive rate along different groups to be equal, that is, 
\begin{equation}
\label{eq.EO}
P(\hat{Y} = 1|A = a, Y = y) = P(\hat{Y} = 1|A = b, Y = y)
\end{equation}
We evaluate our model by multi-class equalized opportunity (Eopp) and equalized odds (Eodd) to quantify the fairness performance \cite{Hardt2016EqualityOO}. The equalized opportunity has two metrics to evaluate the correct prediction. Eopp0 indicates the difference between the two groups' true negative rate (TNR), and Eopp1 illustrates the difference between the two groups' true positive rate (TPR). Equalized odds is the summation of two groups' true positive rate difference and false positive rate difference, which can also address unfair wrong prediction of the model.
% where $y \in \{0, 1\}$, and $a, b \in A$. This metric addresses unfair wrong prediction of the model, and therefore EO will be the suitable fairness metric to evaluate our model. We calculate the degree of EO by calculating the disparity of the $TPR$ and $FPR$ alone different sensitive attributes as follows:
% \begin{equation}
% \label{eq.EO_equation}
%     \sum_{k=1}^{K} |TPR_{k}^{1} - TPR_{k}^{0} + FPR_{k}^{1} - FPR_{k}^{0}|
% \end{equation}
% where $TPR_{k}^{a}$ and $FPR_{k}^{a}$ are the True Positive Rate and the False Positive Rate respectively of target class $k$ and sensitive attribute $a$, this equation can also extend to multi-attribute cases. In conclusion, the optimal EO score becomes $0$ when the True Positive Rate and the False Positive of each target and sensitive class are the same. It indicates that a lower EO represents the prediction is fairer.

% Following most related works, we report EO and accuracy (Acc.) of all experiments in section \ref{sec:result}. 


\subsection{Implementation Details}
We utilize Resnet-18 as the backbone of every state-of-the-art and the baseline CNN used in our experiments. 
In the training phase, the data is augmented by random flipping, rotation, and scaling before being fed into the model. The network is trained for 200 epochs using SGD optimizer \cite{ruder2016overview} with an initial learning rate set as 0.01. We attach three internal classifiers at the end of each residual block. The internal classifier contains one adaptive average pooling layer and a two layer MLP. As introduced in Section \ref{sec:method_main}, the modified loss function in the multi-exit training framework is a weighted summation of the loss from each internal classifier. The coefficients, $\alpha_1$, $\alpha_2$, $\alpha_3$, and $\alpha_f$, are set to $0.3$, $0.45$, $0.6$, and $0.9$, respectively. We set the confidence threshold $\theta=0.85$ for the early exit policy at the inference phase.  

To show that our scheme can be widely applied to most existing frameworks, we conduct experiments on four baselines selected from state-of-art approaches for fairness-aware learning: LNL \cite{kim2019learning}, HSIC \cite{quadrianto2019discovering}, MFD \cite{jung2021fair}, and FSCL \cite{park2022fair}. All the baselines are reproduced by following the recommended hyperparameter settings of the original papers or resources.
 
\begin{table*}[ht]
\centering
\caption{Classification results of the fairness (EO) and accuracy (Acc.) evaluation on the test set of the CelebA dataset. * denotes our own implementation. }
\setlength{\tabcolsep}{4pt}
\label{Table:Compare_SOTA_CelebA}
\begin{tabular}{c||c|c||c|c||c|c||c|c||c|c||c|c}
\hline
\multirow{2}{*}{Methods}           & \multicolumn{2}{c||}{T=a / S=m} &\multicolumn{2}{c||}{T=a / S=y} &\multicolumn{2}{c||}{T=b / S=m}& \multicolumn{2}{c||}{T=b / S=y}& \multicolumn{2}{c||}{T=e / S=m} & \multicolumn{2}{c}{T=e / S=y}  \\ \cline{2-13}
                  & EO            & Acc.          & EO            & Acc.           & EO           & Acc.          & EO           & Acc.          & EO           & Acc.          & EO           & Acc.          \\ \hline \hline
CNN                   & 27.8          & 79.6          & 16.8          & 79.8           & 17.6         & 84.0          & 14.7         & 84.5          & 15.0         & 83.9          & 12.7         & 83.8          \\ 
\textbf{ME-CNN}       & \textbf{23.7} & \textbf{82.3} & \textbf{16.1} & 76.8           & \textbf{12.9} & \textbf{84.8}         & \textbf{12.9} & \textbf{84.8}         & \textbf{10.8} & 82.5          & \textbf{8.1} & 83.6          \\ \hline
LNL                   & 21.8          & 79.9          & 13.7          & 74.3           & 10.7         & 82.3          & 6.8          & 82.3          & 5.0          & 81.6          & 3.3          & 80.3          \\ 
\textbf{ME-LNL}       & \textbf{14.4} & \textbf{82.2} & \textbf{13.1} & 72.7           & \textbf{7.3} & \textbf{82.8} & \textbf{5.5} & \textbf{83.1} & \textbf{2.7} & \textbf{84.0} & \textbf{1.0} & \textbf{82.6} \\ \hline
HSIC*                 & 19.4          & 81.7          & 16.5          & 80.3           & 11.2         & 80.8          & 10.5         & 82.6          & 12.5         & 84.0          & 7.4          & 84.2          \\ 
\textbf{ME-HSIC}      & \textbf{12.9} & 78.8          & \textbf{15.9} & 78.7           & \textbf{3.6} & \textbf{80.9} & \textbf{3.2} & 82.2          & \textbf{7.8} & 83.5          & \textbf{4.1} & 82.0          \\ \hline
FSCL+                 & 6.5           & 79.1          & 12.4          & 79.1           & 4.7          & 82.9          & 4.8          & 84.1          & 3.0          & 83.4          & 1.6          & 83.5          \\ 
\textbf{ME-FSCL+}     & \textbf{5.8}  & 78.2          & \textbf{7.6}  & 76.4           & \textbf{3.6} & 81.2          & \textbf{2.8} & \textbf{84.4} & \textbf{2.0} & \textbf{83.5} & \textbf{1.4} & 80.8          \\ \hline
MFD                   & 7.4           & 78.0          & 14.9          & 80.0           & 7.3          & 78.0          & 5.4          & 78.0          & 8.7          & 79.0          & 5.2          & 78.0          \\ 
\textbf{ME-MFD}       & \textbf{5.8}  & \textbf{78.3} & \textbf{11.4} & 79.5           & \textbf{2.6} & \textbf{82.1} & \textbf{3.3} & \textbf{82.6} & \textbf{1.4} & \textbf{81.9} & \textbf{1.5} & \textbf{84.2} \\ \hline
% FairPrune             &     21.5          &     81.2          &       22.5         &     82.4           &      17.6        &    83.1           &     15.2         &     83.9          &    17.6  &   83.8            &     9.5         &     84.1          \\ 
% \textbf{ME-FairPrune} &               &               &               &                &             &             &              &               &              &               &              &               \\ \hline
\end{tabular}
\end{table*}
% \label{sec:implementation_details}
% \subsubsectionion{ME-CNN}

% \subsubsectionion{\emph{ME-LNL}} 
% We substitute the network {g \circ h} in to ME-CNN with three internal classifiers  and implemented the \emphh{h} with single fully connected layer. For the hyperparameter  



